\begin{frame}{$p^T$ 감쇠: 완주 가능 길이는 어떻게 결정되는가}
  \begin{itemize}
    \item ``완주 확률이 50\%''가 되는 episode 길이:
      \[T=\frac{\log(0.5)}{\log(p)}\]
    \item $p=0.99$이면 \textbf{약 69스텝}, $p=0.995$이면
      \textbf{약 138스텝}.
    \item 1스텝 성공률을 0.5\%p만 올려도 완주 가능한 길이가
      \textbf{두 배}.
  \end{itemize}
  \vskip0.6em
  \begin{block}{읽는 법}
    긴 episode에서는 ``거의 불가능''과 ``가끔 된다''를 가르는 것이
    1스텝 성공률의 \textbf{0.5\%p}다. 12.3절에서 이 표를 5$\times$5
    격자 실험과 대조한다.
  \end{block}
\end{frame}
